\documentclass[11pt,a4paper]{article}

\usepackage[utf8]{inputenc}
\usepackage[T1]{fontenc}
\usepackage[margin=2.5cm]{geometry}
\usepackage{hyperref}
\usepackage{xcolor}
\usepackage{listings}
\usepackage{enumitem}
\usepackage{booktabs}
\usepackage{longtable}
\usepackage{amsmath,amssymb}
\usepackage{fancyhdr}
\usepackage{titlesec}

\definecolor{isls-blue}{HTML}{1a5276}
\definecolor{isls-orange}{HTML}{e67e22}
\definecolor{isls-green}{HTML}{27ae60}
\definecolor{isls-red}{HTML}{c0392b}
\definecolor{isls-gray}{HTML}{7f8c8d}
\definecolor{code-bg}{HTML}{f7f9fb}

\lstdefinestyle{rust}{
  language=C,
  basicstyle=\ttfamily\small,
  keywordstyle=\color{isls-blue}\bfseries,
  commentstyle=\color{isls-gray}\itshape,
  stringstyle=\color{isls-green},
  backgroundcolor=\color{code-bg},
  frame=single,
  rulecolor=\color{isls-gray!30},
  breaklines=true,
  showstringspaces=false,
  tabsize=4,
  morekeywords={fn,let,mut,pub,struct,impl,enum,use,mod,self,Self,
    crate,super,where,trait,for,in,if,else,match,return,async,await,
    Ok,Err,Some,None,Vec,String,HashMap,BTreeMap,Result,Option,
    f64,f32,usize,bool,u32,u64,i64,Box,dyn,Arc},
}
\lstdefinestyle{bash}{
  language=bash,
  basicstyle=\ttfamily\small,
  backgroundcolor=\color{code-bg},
  frame=single,
  rulecolor=\color{isls-gray!30},
  breaklines=true,
  showstringspaces=false,
}

\pagestyle{fancy}
\fancyhf{}
\fancyhead[L]{\small\textcolor{isls-gray}{ISLS I2 Spec --- Infogenetik Evolution}}
\fancyhead[R]{\small\textcolor{isls-gray}{v1.0.0 --- \today}}
\fancyfoot[C]{\thepage}

\titleformat{\section}
  {\Large\bfseries\color{isls-blue}}{\thesection}{1em}{}
\titleformat{\subsection}
  {\large\bfseries\color{isls-blue!80}}{\thesubsection}{1em}{}
\titleformat{\subsubsection}
  {\normalsize\bfseries\color{isls-blue!60}}{\thesubsubsection}{1em}{}

\hypersetup{
  colorlinks=true,
  linkcolor=isls-blue,
  urlcolor=isls-orange,
}

\begin{document}

\begin{titlepage}
\centering
\vspace*{2.5cm}

{\Huge\bfseries\color{isls-blue} ISLS I2 Specification}\\[0.5cm]
{\LARGE\color{isls-orange} Infogenetik Evolution}\\[0.3cm]
{\large\color{isls-gray} Gen-Clustering, Ollama-native Architect,\\
Parallel-Scraping, Codematrix-Feedback}\\[2cm]

{\large
\begin{tabular}{rl}
\textbf{System:}       & ISLS --- Intelligent Semantic Ledger Substrate \\
\textbf{Version:}      & v6.1 (I2) \\
\textbf{Author:}       & Sebastian Klemm \\
\textbf{Date:}         & \today \\
\textbf{Status:}       & Specification \\
\textbf{Depends:}      & I1 (complete), Ollama-Forge-Fix (complete) \\
\textbf{Branch:}       & \texttt{claude/implement-isls-i2-spec} \\
\end{tabular}
}

\vfill

{\small\color{isls-gray}
Mithras Substructure University\\
Repository: \texttt{https://github.com/LashSesh/graphity}\\
License: MIT
}
\end{titlepage}

\tableofcontents
\newpage


% ══════════════════════════════════════════════════════════════════════════════
\section{Executive Summary}
% ══════════════════════════════════════════════════════════════════════════════

I1 installed measurement: 5D Codematrix, Resonites, Mikro/Meso-Gates, 
Fitness-R\"uckkopplung. The system can now \emph{see} what it generates. 
I2 makes it \emph{act} on what it sees.

Three problems to solve:

\begin{enumerate}[nosep]
\item \textbf{The Architect needs OpenAI.} The Chat entity-extraction 
      in the Studio calls GPT-4o for every message. On the VPS with 
      48~GB RAM and Qwen2.5-Coder 32B sitting idle during chat, this 
      is wasteful. I2 makes the Architect Ollama-native: entity 
      extraction via the local model. Zero external API dependency 
      for the full pipeline.
\item \textbf{Norms are flat.} 33 norms (31 builtin + 2 auto) in 
      a flat list. No clustering. No co-activation tracking. No 
      gene structure. I2 introduces Gen-Clustering: norms that 
      consistently co-activate form genes. Genes compose into 
      chromosomes. The Infogenetik becomes operational.
\item \textbf{Scraping is single-threaded.} The \"Olbohrinsel 
      processes one keyword per hour. On an 8-core server with 
      600~Mbit/s, this wastes 87\% of capacity. I2 parallelizes 
      the scraping pipeline and accelerates keyword rotation.
\end{enumerate}

\subsection{What I2 Produces}

\begin{enumerate}[nosep]
\item \textbf{Ollama Architect:} The session chat 
      (\texttt{POST /api/session/\{id\}/message}) uses the local 
      Ollama model for entity extraction when no API key is 
      provided. Full ISLS pipeline without external dependencies.
\item \textbf{Gen-Clustering:} Automatic detection of norm 
      co-activation clusters from \texttt{metrics.jsonl}. 
      Stored as \texttt{genome.json}. CLI: \texttt{isls norms 
      genome}.
\item \textbf{Parallel Scraping:} The mass-scrape endpoint 
      processes multiple keywords simultaneously. 8 workers 
      on 8 cores. 5$\times$ throughput.
\item \textbf{Codematrix in Fitness:} Fitness update uses 
      not just compile-success but the average Codematrix 
      resonance of the generated files. Better code $=$ higher 
      fitness, even if both runs compile.
\end{enumerate}


% ══════════════════════════════════════════════════════════════════════════════
\section{W1: Ollama-Native Architect}
\label{sec:w1}
% ══════════════════════════════════════════════════════════════════════════════

\subsection{Problem}

\texttt{POST /api/session/\{id\}/message} in \texttt{architect.rs} 
calls the OpenAI API for entity extraction. Without an API key, it 
falls back to ``manual mode'' --- the user must add entities via 
REST. This defeats the purpose of the chat interface.

\subsection{Solution}

The Architect uses the same Oracle abstraction as the Forge. When 
\texttt{OracleConfig} has \texttt{use\_ollama = true}, the 
Architect constructs an \texttt{OllamaOracle} and calls it for 
entity extraction.

\subsection{Prompt Adaptation}

The Architect prompt currently assumes GPT-4o's instruction-following 
capability. For Qwen2.5-Coder 32B, the prompt needs to be tighter:

\begin{lstlisting}[style=rust]
pub fn build_architect_prompt_ollama(
    history: &[SessionMessage],
    current_spec: &AppSpec,
    user_message: &str,
) -> String {
    format!(
        r#"You are a software architect. The user describes an 
application. Extract entities with their fields.

CURRENT STATE:
App name: {app_name}
Entities: {entities}

USER MESSAGE: {message}

Respond with ONLY a JSON object. No markdown. No explanation.
Example format:
{{"message":"Created Trade entity","updates":{{"entities_add":[{{"name":"Trade","fields":[{{"name":"price","field_type":"f64"}}]}}],"entities_modify":[],"entities_remove":[],"infra":{{}}}}}}

JSON:"#,
        app_name = current_spec.app_name,
        entities = format_entities_compact(current_spec),
        message = user_message,
    )
}
\end{lstlisting}

Key differences from the GPT-4o prompt:
\begin{itemize}[nosep]
\item Shorter, more direct. Local models follow shorter prompts 
      better.
\item Explicit JSON-only instruction with example.
\item No conversation history injection (reduces context length; 
      current spec is sufficient context).
\end{itemize}

\subsection{Fallback Chain}

\begin{enumerate}[nosep]
\item Session has \texttt{api\_key}: use OpenAI (best quality).
\item Server has \texttt{oracle\_config.api\_key}: use OpenAI.
\item Server has \texttt{use\_ollama = true}: use Ollama 
      (local, free, slightly lower quality).
\item Nothing: manual mode (existing behavior, no change).
\end{enumerate}

\subsection{JSON Parsing Robustness}

Local models sometimes wrap JSON in markdown fences or add 
preamble text. The parser must:
\begin{enumerate}[nosep]
\item Try parsing the raw response as JSON.
\item If that fails: extract the first \texttt{\{...\}} block 
      using brace-matching.
\item If that fails: try removing \texttt{```json} and 
      \texttt{```} wrappers.
\item If all fail: treat as manual-mode message (log warning, 
      don't crash).
\end{enumerate}

\subsection{Modified Files}

\begin{longtable}{lp{10cm}}
\toprule
\textbf{File} & \textbf{Change} \\
\midrule
\texttt{isls-gateway/src/architect.rs}
  & Add \texttt{build\_architect\_prompt\_ollama()}, 
    Oracle-based entity extraction, JSON parsing 
    robustness. \\
\texttt{isls-gateway/src/session.rs}
  & Use \texttt{OracleConfig} from AppState to construct 
    Oracle for architect calls. \\
\texttt{isls-gateway/src/lib.rs}
  & Pass \texttt{OracleConfig} to session/architect 
    handlers. \\
\bottomrule
\end{longtable}

\subsection{W1 Verification}

\begin{enumerate}[nosep]
\item Start \texttt{isls serve --port 8420 --ollama} (no API key).
\item Open Studio, type ``Pet shop with animals and owners''.
\item System responds with entities (via Ollama, not OpenAI).
\item Entity cards appear in chat.
\item Readiness ring updates.
\item Type ``forge'' --- forge runs with Ollama.
\item \textbf{Full pipeline without any external API.}
\end{enumerate}


% ══════════════════════════════════════════════════════════════════════════════
\section{W2: Gen-Clustering}
\label{sec:w2}
% ══════════════════════════════════════════════════════════════════════════════

\subsection{Problem}

Norms are a flat list. The system doesn't know that CRUD-Entity, 
JWT-Auth, Pagination, and ErrorSystem \emph{always} activate 
together. This co-activation pattern is a ``gene'' --- a cluster 
of norms that form a functional unit.

\subsection{Data Source}

\texttt{metrics.jsonl} contains \texttt{norms\_activated} for 
every generation. After 20+ entries, co-activation patterns 
are detectable.

\subsection{Algorithm}

\begin{lstlisting}[style=rust]
pub fn cluster_genes(
    metrics: &[GenerationMetrics],
    min_coactivation: f64,  // default 0.8
) -> Vec<Gene> {
    // 1. Build co-activation matrix
    let all_norms: BTreeSet<String> = metrics.iter()
        .flat_map(|m| m.norms_activated.iter().cloned())
        .collect();
    
    let n = all_norms.len();
    let norm_list: Vec<&String> = all_norms.iter().collect();
    let mut coact = vec![vec![0u32; n]; n];
    let mut counts = vec![0u32; n];
    
    for m in metrics {
        let active: Vec<usize> = norm_list.iter().enumerate()
            .filter(|(_, name)| m.norms_activated.contains(*name))
            .map(|(i, _)| i)
            .collect();
        for &i in &active {
            counts[i] += 1;
            for &j in &active {
                coact[i][j] += 1;
            }
        }
    }
    
    // 2. Compute Jaccard similarity
    let mut sim = vec![vec![0.0f64; n]; n];
    for i in 0..n {
        for j in 0..n {
            if counts[i] + counts[j] > coact[i][j] {
                sim[i][j] = coact[i][j] as f64 
                    / (counts[i] + counts[j] - coact[i][j]) as f64;
            }
        }
    }
    
    // 3. Agglomerative clustering
    //    Merge norms with sim >= min_coactivation
    //    into gene clusters.
    agglomerative_cluster(&sim, &norm_list, min_coactivation)
}
\end{lstlisting}

\subsection{Gene Data Structure}

\begin{lstlisting}[style=rust]
#[derive(Debug, Clone, Serialize, Deserialize)]
pub struct Gene {
    pub id: String,           // "GENE-0001"
    pub name: String,         // auto-generated or user-named
    pub norms: Vec<String>,   // norm IDs in this cluster
    pub coactivation: f64,    // avg pairwise coactivation
    pub fitness: f64,         // avg fitness of member norms
    pub activation_count: u64,
    pub domains: Vec<String>, // domains where this gene activated
}

#[derive(Debug, Clone, Serialize, Deserialize)]
pub struct Genome {
    pub genes: Vec<Gene>,
    pub generation: u64,      // how often genome was recomputed
    pub last_updated: String,
    pub total_metrics: usize, // metrics entries used
}
\end{lstlisting}

\subsection{Storage}

\texttt{\textasciitilde/.isls/genome.json}. Recomputed whenever 
\texttt{isls norms genome} is called or when 10 new metrics 
entries have accumulated since last computation.

\subsection{Auto-Naming}

Genes are named by their dominant norm:
\begin{itemize}[nosep]
\item Gene containing CRUD + Auth + Pagination $\to$ ``WebApp-Core''
\item Gene containing Docker + Nginx + Env $\to$ ``Deployment''
\item Gene containing only one norm $\to$ norm's own name (singleton)
\end{itemize}

\subsection{CLI}

\begin{lstlisting}[style=bash]
isls norms genome

ISLS Genome (Generation 3, 19 metrics entries)
---
GENE-0001  "WebApp-Core"  coact=1.00  fitness=0.95
  ISLS-NORM-0042 (CRUD-Entity)
  ISLS-NORM-0088 (JWT-Auth)
  ISLS-NORM-0096 (Pagination)
  ISLS-NORM-0103 (ErrorSystem)
  Domains: pet-shop, hotel, clinic, gym, ... (15)

GENE-0002  "Infra-Web"  coact=0.95  fitness=0.92
  ISLS-NORM-INFRA-WEB
  ISLS-NORM-INFRA-DB
  ISLS-NORM-INFRA-AUTH
  ISLS-NORM-INFRA-FRONTEND
  ISLS-NORM-INFRA-DOCKER
  Domains: all

Singletons (not yet clustered):
  ISLS-NORM-INFRA-CLI (2 activations)
  ISLS-NORM-INFRA-LIB (0 activations)
\end{lstlisting}

\subsection{New Files}

\begin{longtable}{lp{10cm}}
\toprule
\textbf{File} & \textbf{Purpose} \\
\midrule
\texttt{isls-norms/src/genome.rs}
  & \texttt{Gene}, \texttt{Genome}, \texttt{cluster\_genes()}, 
    persistence, CLI output. \\
\bottomrule
\end{longtable}

\subsection{Modified Files}

\begin{longtable}{lp{10cm}}
\toprule
\textbf{File} & \textbf{Change} \\
\midrule
\texttt{isls-norms/src/lib.rs}
  & \texttt{pub mod genome;} \\
\texttt{isls-cli/src/cmd\_norms.rs}
  & Add \texttt{isls norms genome} subcommand. \\
\texttt{isls-cli/src/main.rs}
  & Route \texttt{genome} subcommand. \\
\bottomrule
\end{longtable}

\subsection{W2 Verification}

\begin{enumerate}[nosep]
\item Need 20+ entries in \texttt{metrics.jsonl} (from fulltest 
      or Studio usage).
\item \texttt{isls norms genome} shows gene clusters.
\item WebApp-Core gene (CRUD + Auth + Pagination + Error) detected.
\item Singletons listed separately.
\end{enumerate}


% ══════════════════════════════════════════════════════════════════════════════
\section{W3: Parallel Scraping}
\label{sec:w3}
% ══════════════════════════════════════════════════════════════════════════════

\subsection{Problem}

The current mass-scrape endpoint processes keywords sequentially. 
One keyword $\to$ search $\to$ 5 repos $\to$ scrape each $\to$ 
next keyword. With GitHub rate limiting (2s between API calls), 
this is slow. But the rate limit is only on \emph{GitHub search}. 
The actual scraping (clone + parse + observe) is CPU-bound and 
can be parallelized.

\subsection{Solution}

Separate the pipeline into two phases:

\begin{enumerate}[nosep]
\item \textbf{Search phase (sequential):} Query GitHub API for 
      all keywords. Collect all repo URLs. Respect rate limit 
      (2s between searches). This takes $N_\text{keywords} \times 2$ 
      seconds.
\item \textbf{Scrape phase (parallel):} Scrape all collected repos 
      in parallel using a worker pool. $W$ workers (default: 
      \texttt{num\_cpus / 2}, max 4). Each worker: clone $\to$ 
      parse $\to$ observe $\to$ delete temp dir.
\end{enumerate}

\subsection{Implementation}

\begin{lstlisting}[style=rust]
async fn mass_scrape_parallel(
    keywords: Vec<String>,
    results_per_keyword: usize,
    state: AppState,
) {
    // Phase 1: Collect all repos (sequential, rate-limited)
    let mut all_repos: Vec<RepoInfo> = Vec::new();
    for keyword in &keywords {
        let repos = github_search(keyword, results_per_keyword).await;
        all_repos.extend(repos);
        publish_event(&state, DiscoverKeyword { ... });
        tokio::time::sleep(Duration::from_secs(2)).await;
    }
    
    // Phase 2: Scrape all repos (parallel)
    let semaphore = Arc::new(Semaphore::new(4)); // 4 workers
    let mut handles = Vec::new();
    
    for repo in all_repos {
        let sem = semaphore.clone();
        let state = state.clone();
        handles.push(tokio::spawn(async move {
            let _permit = sem.acquire().await.unwrap();
            let result = scrape_repo(&repo, &state).await;
            publish_event(&state, DiscoverRepo { ... });
            result
        }));
    }
    
    let results = futures::future::join_all(handles).await;
    publish_event(&state, DiscoverComplete { ... });
}
\end{lstlisting}

\subsection{Hourly Scrape Enhancement}

Update \texttt{scrape\_hourly.sh} to send 5 keywords per run 
instead of 1:

\begin{lstlisting}[style=bash]
#!/bin/bash
KEYWORDS=$(shuf -n 5 /home/mithra/isls/scrape_keywords.txt \
  | paste -sd ',' -)
# Convert to JSON array
JSON_KW=$(echo "$KEYWORDS" | sed 's/,/","/g' | sed 's/^/"/' | sed 's/$/"/')
curl -s -X POST http://localhost:8420/api/discover/mass-scrape \
  -H "Content-Type: application/json" \
  -d "{\"keywords\":[$JSON_KW],\"results_per_keyword\":3}" \
  >> /home/mithra/isls/scrape.log 2>&1
\end{lstlisting}

Result: 5 keywords $\times$ 3 repos $\times$ 24 hours = 
\textbf{360 repos/day} (up from 120).

\subsection{Modified Files}

\begin{longtable}{lp{10cm}}
\toprule
\textbf{File} & \textbf{Change} \\
\midrule
\texttt{isls-gateway/src/discover.rs}
  & Refactor mass-scrape to two-phase parallel pipeline. 
    Add semaphore-based worker pool. \\
\bottomrule
\end{longtable}

\subsection{W3 Verification}

\begin{enumerate}[nosep]
\item \texttt{POST /api/discover/mass-scrape} with 5 keywords 
      runs faster than sequential.
\item WebSocket events still arrive in order per keyword.
\item No GitHub rate limit violations (search phase is sequential).
\item 4 concurrent scrape operations visible in server logs.
\end{enumerate}


% ══════════════════════════════════════════════════════════════════════════════
\section{W4: Codematrix-Enhanced Fitness}
\label{sec:w4}
% ══════════════════════════════════════════════════════════════════════════════

\subsection{Problem}

I1 Fitness uses binary success ($r_t \in \{0, 1\}$). Two runs 
that both compile get the same fitness boost, even if one 
produced clean code (Codematrix resonance 0.85) and the other 
produced messy code (resonance 0.45).

\subsection{Solution}

Replace binary $r_t$ with continuous Codematrix resonance:

\begin{equation}
r_t = \begin{cases}
\bar{R}_\text{codematrix} & \text{if compile success} \\
0 & \text{if compile failed}
\end{cases}
\end{equation}

where $\bar{R}_\text{codematrix}$ is the average resonance 
product across all generated files in this run 
(from \texttt{ForgeStats.codematrix\_avg}).

\begin{equation}
\phi_{t+1}(n) = \alpha \cdot \phi_t(n) + (1 - \alpha) \cdot r_t
\end{equation}

with $\alpha = 0.9$ (unchanged).

\subsection{Effect}

\begin{itemize}[nosep]
\item A norm that consistently produces high-resonance code 
      gets $\phi \to 0.85+$.
\item A norm that compiles but produces messy code gets 
      $\phi \to 0.5$.
\item A norm that fails to compile gets $\phi \to 0$.
\item \textbf{The system learns quality, not just success.}
\end{itemize}

\subsection{Modified Files}

\begin{longtable}{lp{10cm}}
\toprule
\textbf{File} & \textbf{Change} \\
\midrule
\texttt{isls-norms/src/fitness.rs}
  & \texttt{update\_fitness()} takes \texttt{Option<f64>} 
    codematrix resonance instead of \texttt{bool} success. \\
\texttt{isls-forge-llm/src/staged\_closure.rs}
  & Compute average codematrix resonance from generated files. 
    Pass to \texttt{update\_fitness()}. \\
\texttt{isls-cli/src/main.rs}
  & Pass codematrix data through to fitness update in CLI 
    forge paths. \\
\bottomrule
\end{longtable}

\subsection{W4 Verification}

\begin{enumerate}[nosep]
\item Run two forge generations.
\item \texttt{isls norms fitness} shows fitness values between 
      0 and 1 (not just 0.5 or 0.9x).
\item Fitness reflects code quality, not just compile success.
\end{enumerate}


% ══════════════════════════════════════════════════════════════════════════════
\section{Dependency Graph}
% ══════════════════════════════════════════════════════════════════════════════

\begin{verbatim}
W1 (Ollama Architect)     W3 (Parallel Scraping)
         |                        |
         v                        v
W4 (Codematrix Fitness)   W2 (Gen-Clustering)
\end{verbatim}

W1 and W3 are independent. W4 depends on W1 (needs Ollama 
forge to produce real codematrix data). W2 depends on 
metrics.jsonl having entries (from W1 or from previous runs).

Recommended: W1 $\to$ W3 $\to$ W4 $\to$ W2.


% ══════════════════════════════════════════════════════════════════════════════
\section{Constraints for Claude Code}
% ══════════════════════════════════════════════════════════════════════════════

\begin{enumerate}
\item \textbf{Ollama Architect} reuses the existing \texttt{Oracle} 
      trait. Construct \texttt{OllamaOracle} from 
      \texttt{AppState.oracle\_config}. No new HTTP client.
\item \textbf{JSON parsing robustness} is critical. Local models 
      produce inconsistent formatting. The parser must not crash 
      on unexpected output.
\item \textbf{Gen-Clustering} reads \texttt{metrics.jsonl}. If 
      fewer than 10 entries exist, return empty genome with 
      message ``Not enough data.''
\item \textbf{Parallel scraping} uses \texttt{tokio::sync::Semaphore} 
      with max 4 permits. Do NOT spawn unbounded tasks.
\item \textbf{Codematrix fitness} must be backward compatible: if 
      codematrix data is not available (old metrics entries), 
      fall back to binary success.
\item \textbf{Do not modify} isls-merkaba, isls-reader, 
      isls-code-topo, isls-hypercube, isls-types, isls-chat.
\item All existing tests MUST pass.
\item Branch: \texttt{claude/implement-isls-i2-spec}.
\item Commit format: \texttt{I2/W\{n\}: <description>}.
\end{enumerate}


% ══════════════════════════════════════════════════════════════════════════════
\section{Acceptance Criteria}
% ══════════════════════════════════════════════════════════════════════════════

\begin{longtable}{clp{9cm}}
\toprule
\textbf{\#} & \textbf{Pass} & \textbf{Criterion} \\
\midrule
1 & & \texttt{cargo build --workspace} compiles. \\
2 & & \texttt{cargo test --workspace} passes. \\
3 & & Studio chat works with \texttt{--ollama} only (no API key). 
      Entities extracted, readiness computed. \\
4 & & Forge from Studio uses Ollama. Real LLM calls visible 
      (duration $> 30$s, tokens $> 0$). \\
5 & & \texttt{isls norms genome} shows gene clusters (needs 
      10+ metrics entries). \\
6 & & Mass-scrape with 5 keywords uses parallel workers. \\
7 & & Fitness values reflect codematrix resonance, not just 
      binary success. \\
8 & & Full pipeline (chat $\to$ forge $\to$ compile) works 
      without any external API. \\
\bottomrule
\end{longtable}

I2 is passed when criteria 1--4 and 8 are fulfilled. 
Criteria 5 requires accumulated metrics. Criteria 6--7 
are improvements that can be verified independently.


% ══════════════════════════════════════════════════════════════════════════════
\section{After I2: Full Sovereignty}
% ══════════════════════════════════════════════════════════════════════════════

With I2 complete, ISLS runs entirely on the VPS:

\begin{itemize}[nosep]
\item \textbf{Chat:} Ollama (free).
\item \textbf{Forge:} Ollama (free).
\item \textbf{Scraping:} GitHub API (free, rate-limited).
\item \textbf{Norms:} Local (\texttt{norms.json}).
\item \textbf{Fitness:} Local (\texttt{fitness.json}).
\item \textbf{Genome:} Local (\texttt{genome.json}).
\item \textbf{Metrics:} Local (\texttt{metrics.jsonl}).
\item \textbf{Studio:} Self-hosted (port 8420).
\end{itemize}

Zero external dependencies. Zero recurring costs beyond 
VPS hosting. The system learns, generates, and improves 
autonomously. The \"Olbohrinsel pumps. The Genome crystallizes. 
The Fitness converges. The operator sleeps.

\vfill
\begin{center}
\color{isls-gray}\rule{0.5\textwidth}{0.4pt}\\[0.5em]
{\small End of I2 Specification.\\[0.3em]
Souver\"anit\"at ist nicht Unabh\"angigkeit von Technologie.\\
Souver\"anit\"at ist Unabh\"angigkeit von Erlaubnis.}
\end{center}

\end{document}
